\section{Inference: why and how the placebo nulls work}
\label{sec:temoin}

This is the most important section for understanding the results of the
journal version. It explains why estimated entropy measures are biased and
move on their own, how a resampling null corrects for it, why a second,
volatility-preserving null is needed for claims about dependence, and how the
size and power of the test are checked.

\subsection{The bias of a convex function of estimated moments}

\begin{lemme}[Jensen applied to an estimator]
If $\hat\theta$ is an (approximately) unbiased estimator of $\theta$ and $f$
is convex, then $\E[f(\hat\theta)]\ge f(\theta)$.
\end{lemme}

\noindent $C=-\tfrac12\log\det R$ is a convex function of $R$ ($-\log\det$ is
convex on positive definite matrices), and $\DD$ is non-negative and zero only
at the Gaussian point. Both are therefore \textbf{biased upward}: estimation
noise looks like information.

\subsection{How large is the bias of coupling? (Proposition 3 of the paper)}

Take $N$ \emph{independent} Gaussian assets, so the true coupling is $C=0$,
and $T$ observations. The likelihood-ratio statistic for testing $H_0:R=I$
(``sphericity'' of the correlation matrix) is
\[
\Lambda=-T\log\det\hat R .
\]
Classical asymptotic theory says that $-2\log(\text{likelihood ratio})$ is
approximately $\chi^2$ with as many degrees of freedom as free parameters
restricted by $H_0$. Here $H_0$ fixes the $N(N-1)/2$ off-diagonal
correlations, so $\Lambda\approx\chi^2_{N(N-1)/2}$. Bartlett (1950) showed that
replacing $T$ by $T'=T-1-(2N+5)/6$ makes the approximation much better in
small samples. Since a $\chi^2_k$ variable has mean $k$,
\[
\E\bigl[-\log\det\hat R\bigr]\approx\frac{N(N-1)}{2T'}
\quad\Longrightarrow\quad
\E\bigl[C(\hat R)\bigr]\approx b(N,T)=\frac{N(N-1)}{4\bigl(T-1-(2N+5)/6\bigr)} .
\]

\begin{exemple}[The paper's panels]
For $N=20$, $T=252$: $T'=251-7.5=243.5$ and $b=380/974=\BIASTWENTY$~nat. For
$N=48$, $T=252$: $T'=251-16.8=234.2$ and $b=2256/936.8=\BIASFORTYEIGHT$~nats.
The simulation of the paper (table~2) gives \BIASTWENTYMC~nat for $N=20$: the
formula is accurate. The bias is as large as a third of the effect we want
to measure on stocks ($0.99$ nat).
\end{exemple}

\begin{intuition}
Why fourth moments barely matter: for two independent variables with finite
fourth moments, $\sqrt T\hat\rho\to\N(0,1)$ whatever their marginal laws,
because $\operatorname{Var}(XY)=\operatorname{Var}(X)\operatorname{Var}(Y)$
when $X$ and $Y$ are independent. The bias formula only uses the variance of
$\hat\rho$, so it holds for Student-$t_5$ data too (the simulation gives
\BIASTWENTYT).
\end{intuition}

\paragraph{Two practical consequences.} (1) Coupling levels on panels of
different size are not comparable: subtract $b(N,T)$ first. The
bias-corrected effective number of modes is $N\exp\{-2[C(\hat R)-b]/N\}$.
(2) A \emph{difference} between two windows of the same length (stress minus
calm) has the bias removed to first order. This is why every test in the
paper is on a gap, not on a level.

\begin{piege}
Shrinkage (Ledoit--Wolf) pulls $\hat R$ towards $I$ and therefore removes most
of the bias (\BIASTWENTYLW~nat in the simulation). It does not remove the
\emph{dynamic} artefact of the next subsection.
\end{piege}

\subsection{The mechanical sawtooth}

When an extreme day enters the window, the estimated moments jump, and a
convex function of these moments jumps \emph{upward whatever the sign of the
day}. As the day ages (exponential weighting) or leaves the window, the index
decays slowly.

\paragraph{A short calculation.} Take a window of $T$ standardised returns
whose sample correlation is exactly $I$ (independent assets), and add one day
on which every asset moves by the same amount $a$ standard deviations. Ignoring
the change in the means, the new second-moment matrix is $I+x\mathbf 1\mathbf
1^\top$ with $x=a^2/T$, and the new correlation matrix is
$R'=(I+x\mathbf 1\mathbf 1^\top)/(1+x)$. By the matrix determinant lemma,
$\det(I+x\mathbf 1\mathbf 1^\top)=1+Nx$, so
\[
\Delta C=-\tfrac12\log\det R'=\tfrac12\bigl[N\log(1+x)-\log(1+Nx)\bigr]
\approx\frac{N(N-1)}{4}\,x^2\quad(x\text{ small}).
\]
The first-order terms cancel and what remains is positive: a single common day
raises coupling, whatever its sign. For $N=20$, $T=252$ and a 10-standard-deviation
day ($x=0.40$), $\Delta C=\tfrac12[20\log1.40-\log8.94]=2.2$~nats. The day then
leaves the estimate only at the rate at which the window forgets: with
exponential weights of half-life 60 days, its weight halves every 60 days.

On i.i.d.\ Gaussian data the paper's pipeline finds $37$ positive jumps of
$\JJ$, no negative one, and an apparent half-life of $35$ days, with no event
at all.

\subsection{The resampling test and why its $p$-value is valid}

\paragraph{Construction.} Let $S(X)$ be a statistic computed by the
\emph{whole} pipeline on the $T\times N$ panel $X$ (rolling estimation, stress
regime, jump detection, gap). A null is a way of generating panels $X^*$ that
keep some features of $X$ and destroy others. We draw $B=200$ of them and
compute $S(X^*_1),\dots,S(X^*_B)$.

\paragraph{The $p$-value.} In the predicted direction,
\[
\hat p=\frac{1+\#\{b:\ S(X^*_b)\ge S(X)\}}{B+1}.
\]
If, under the null hypothesis, the observed panel is just one more draw from
the same mechanism, then $S(X),S(X^*_1),\dots,S(X^*_B)$ are
\emph{exchangeable}: the observed value is equally likely to have any rank
among the $B+1$ values. The probability that it ranks in the top $k$ is
$k/(B+1)$, so $P(\hat p\le\alpha)\le\alpha$ exactly, without any asymptotic
argument. The ``$+1$'' is what makes this exact; with $B=200$ the smallest
possible $p$-value is $1/201\approx0.005$.

\begin{piege}
Reporting $z=(S(X)-\bar S^*)/s_{S^*}$ (``six standard deviations above the
null'') is a useful effect size, but it is not a $p$-value: the null
distribution need not be normal. The paper reports $\hat p$ for inference and
$z$ only as a description.
\end{piege}

\subsection{Null 1: i.i.d.\ resampling of days}

Draw rows of $X$ with replacement. Each row (one day, $N$ returns) is intact,
so the null keeps each asset's fat tails, each day's cross-sectional
correlation and common extreme days. It destroys \emph{all} dynamics: there
is no volatility clustering, no persistence, no crisis that lasts. It answers
the question ``beyond the fact that turbulent periods contain big days, does
the measure detect any dynamics?''

\subsection{Null 2: constant-correlation GARCH (filtered historical simulation)}

\paragraph{Why a second null.} Under the i.i.d.\ null, a ``stress'' period is
a random cluster of big days. In real data, stress comes with volatility
clustering: high-volatility days follow each other for weeks. Could
volatility clustering alone, with correlations that never change, produce the
coupling gap? Two mechanisms make this plausible. Forbes and Rigobon (2002)
showed that measured correlation rises mechanically when the variance of a
common factor rises. And our stress regime is defined by high volatility, so
it selects precisely the windows where this happens.

\paragraph{Construction.} For each asset $i$, fit a GJR-GARCH(1,1):
\[
r_{i,t}=\mu_i+\varepsilon_{i,t},\qquad \varepsilon_{i,t}=\sqrt{h_{i,t}}\,z_{i,t},\qquad
h_{i,t}=\omega_i+\bigl(\alpha_i+\gamma_i\mathbf 1\{\varepsilon_{i,t-1}<0\}\bigr)\varepsilon_{i,t-1}^2+\beta_ih_{i,t-1}.
\]
The term $\gamma_i$ makes volatility react more to falls than to rises (the
leverage effect). Estimate by maximising the Gaussian likelihood (quasi-MLE,
consistent even if $z$ is not Gaussian). Compute the standardised residuals
$\hat z_{i,t}=\hat\varepsilon_{i,t}/\sqrt{\hat h_{i,t}}$ and store them as a
$T\times N$ matrix $\hat Z$.

To simulate: draw rows of $\hat Z$ with replacement (as in null 1), then
rebuild each asset's returns by running its own GARCH recursion forward with
the drawn shocks. The result has three properties:
\begin{enumerate}
\item each asset keeps its own volatility clustering and leverage effect;
\item all assets receive the shocks of the same historical day, so a common
large day still hits everyone, and volatility bursts are common;
\item the correlation of the shocks $z_t$ is the same every day: the
\emph{conditional} correlation is constant (Bollerslev's CCC model).
\end{enumerate}
So volatility moves as in a GARCH world and correlations do not. A coupling
gap beyond this null is a genuine dependence dynamic.

\begin{intuition}
Null 1 says: ``it is not just big days''. Null 2 says: ``it is not just
volatility''. The 21-day block null that the earlier version rejected kept
\emph{both} volatility and correlation dynamics over a month, i.e.\ the crisis
itself. Null 2 keeps the volatility part only, which is exactly the
counterfactual the coupling claim needs.
\end{intuition}

\subsection{A trend is not a stress effect: the within-year gap}

Correlations between stocks rose over 1990--2022, and stress dates cluster in
a few years (2008, 2020). If coupling simply trends up and stress dates happen
to be late, the stress-minus-calm gap picks up the trend. The fix is a
regression with year fixed effects:
\[
C_t=\theta\,\mathbf 1\{t\in\text{stress}\}+\sum_y\delta_y\mathbf 1\{t\in\text{year }y\}+u_t .
\]
By the Frisch--Waugh--Lovell theorem, $\hat\theta$ is the slope of the
regression of the within-year demeaned $C_t$ on the within-year demeaned
stress dummy: it compares stress and calm dates \emph{of the same year}. The
paper finds that $\hat\theta$ is about \FESHRINK\ smaller than the raw gap, but
still beyond both nulls.

\subsection{Checking size and power by Monte Carlo}

A test has the right \emph{size} if, when the null hypothesis is true, it
rejects in about $5\%$ of samples at the $5\%$ level; its \emph{power} is the
rejection rate when the hypothesis is false. To check both, simulate panels
from known processes, apply the whole test to each, and count rejections:
\begin{itemize}
\item i.i.d.\ Student-$t$, constant correlation: no dynamics at all. Both
nulls should reject about 5\%.
\item GARCH with constant correlation: volatility clustering, no
correlation dynamics. A good test of \emph{dependence} should still reject
about 5\%. The i.i.d.\ null rejects \MCGARCHIID\ of the time, the GARCH null
\MCGARCHFHS.
\item Markov-switching volatility, constant correlation: same question with
regime-type volatility.
\item Markov-switching volatility \emph{and} correlation (0.30 calm, 0.50 in
stress): the effect is real; rejection rates measure power
(\MCPOWFHS\ for the GARCH null).
\end{itemize}

\subsection{Testing many hypotheses at once: Holm's procedure}

The verdict table tests $m$ statistics. If each is tested at 5\%, the chance
of at least one false rejection grows quickly with $m$. The Holm (1979)
procedure controls this \emph{family-wise error rate}. Sort the $p$-values
$p_{(1)}\le\dots\le p_{(m)}$; reject $H_{(1)}$ if $p_{(1)}\le\alpha/m$, then
$H_{(2)}$ if $p_{(2)}\le\alpha/(m-1)$, and so on, stopping at the first
failure. Equivalently, the adjusted $p$-value of the $k$-th smallest is
$\max_{j\le k}\min\{1,(m-j+1)p_{(j)}\}$. Holm is always at least as powerful
as Bonferroni (which uses $\alpha/m$ for all) and needs no assumption on the
dependence between tests.

\subsection{Comparing nested forecasts: Clark--West}

H7 compares a small model (volatility only, forecast $f_1$) with a larger one
that adds an entropy regressor (forecast $f_2$). Under the null that the extra
regressor is useless, its true coefficient is zero, but the larger model
\emph{estimates} it, which adds noise: $f_2$ has a slightly larger mean squared
error even when the regressor is irrelevant. The Diebold--Mariano test on
$d_t=(y_t-f_{1,t})^2-(y_t-f_{2,t})^2$ is therefore biased against the larger
model and under-rejects. Clark and West (2007) add back the noise term:
\[
\tilde d_t=(y_t-f_{1,t})^2-\bigl[(y_t-f_{2,t})^2-(f_{1,t}-f_{2,t})^2\bigr],
\]
and test $\E[\tilde d_t]=0$ against $\E[\tilde d_t]>0$ with a HAC standard
error. Why this works: if the extra regressor is useless, $f_2$ equals $f_1$
plus pure estimation noise $\nu_t$ uncorrelated with the forecast error, so
$\E[(y-f_2)^2]=\E[(y-f_1)^2]+\E[\nu^2]$. The DM differential then has mean
$-\E[\nu^2]<0$, and $(f_1-f_2)^2=\nu^2$ is exactly the term that re-centres it
at zero. The test is one-sided and approximately normal.

\begin{papier}
Verdicts (tables 4 and 5 of the journal version): coupling under stress
passes both nulls on both panels, also within years; crisis composition (H8),
the range and daily half-life of $\JJ$ and the largest episode pass the
i.i.d.\ null. One-sided jumps, the shape deficit under stress, the
tails--correlation link (H6) and the daily sign (H5) are produced just as
strongly, or more, by the null.
\end{papier}
